% results_interpretation.tex
\subsection{Interpretation Tools}

After \citet{wulff-mata-2025}, \texttt{sfa\_anchor()} scores each item's
cosine to every construct centroid or construct-name embedding, flagging
items weak or beaten on their own construct:

{\scriptsize
\begin{verbatim}
> anc <- sfa_anchor(fit, anchor = "centroid")
> anc
Semantic anchoring
  Method: Wulff & Mata (2025)
  Anchor: centroid
  Constructs: 5 | Items: 50

[50 x 5 item-by-construct cosine matrix omitted]

  Weakest own-construct similarity (review/removal candidates):
    N12          Neuroticism        own=0.69   <-- higher on Extraversion (0.72)
        "I am relaxed most of the time."
    O44          Openness           own=0.69
        "I am not interested in abstract ideas."
    A23          Agreeableness      own=0.71   <-- higher on Extraversion (0.71)
        "I insult people."
    A25          Agreeableness      own=0.71
        "I am not interested in other people's problems."
    [list abridged]
\end{verbatim}
}

Centroid anchoring mostly flags reverse-keyed items worded nearer another
domain---N12 sits nearer the Extraversion centroid (.72) than its own
(.69); label anchoring (same Qwen3-Embedding-8B space) flagged an
overlapping set, A25 lowest (own .36 vs.\ .37 on Neuroticism), then E6,
A27, C36, N12, and C32.

\texttt{sfa\_project()} applies the semantic projection of
\citet{grand-etal-2022}, placing items on bipolar axes via antonym-pole
difference vectors from the same 8B space:

{\scriptsize
\begin{verbatim}
> prj <- sfa_project(fit,
+                    axes = list(stability   = c(low = "anxious",  high = "calm"),
+                                sociability = c(low = "solitary", high = "sociable")),
+                    pole_embeddings = poles)
> prj
Semantic projection onto 2 axis/axes
  Method: Grand et al. (2022)
  Scale: 0 = low pole, 1 = high pole

Axis 'stability'  [anxious <-> calm]
  lowest: N13=0.28  N11=0.33  N16=0.37  N19=0.38  N15=0.42
  highest: N12=0.82  A30=0.72  E3=0.65  E4=0.64  C37=0.62

Axis 'sociability'  [solitary <-> sociable]
  lowest: O46=0.37  A27=0.38  A21=0.38  N20=0.39  E4=0.41
  highest: E3=0.82  E7=0.81  A22=0.68  E5=0.68  A30=0.65
\end{verbatim}
}

The projections track content, not scoring key: Neuroticism items fill the
anxious pole, yet the most calm-pole item (.82) is again reverse-keyed
N12, while party- and conversation-themed E3 and E7, and other-oriented
A22 and A30, top the sociable pole.

Versus the theoretical Big Five assignment, \texttt{sfa\_congruence()}
scored agreement as moderate by NMI and weak by the chance-corrected ARI
\citep{hubert-arabie-1985}:

{\scriptsize
\begin{verbatim}
> cong_theory <- sfa_congruence(fit, target = big5$factors,
+                               metrics = c("nmi", "ari"))
> cong_theory
Factor structure congruence

  NMI:            0.527 (moderate - higher is better)
  ARI:            0.402 (weak - higher is better)
\end{verbatim}
}

\texttt{sfa\_jinglejangle()} compares label--label with content--content
similarities \citep{wulff-mata-2025, wulff-mata-2026}, here treating the
five 10-item subscales as separate scales (8B item and label vectors):

{\scriptsize
\begin{verbatim}
> jj <- sfa_jinglejangle(sub_items, item_embeddings = sub_emb,
+                        label_embeddings = jj_lab_emb)
> jj
Jingle-jangle detection across 5 scales
  Method: Wulff & Mata (2025, 2026)
  Flag threshold |content - label| >= 0.20

  No jingle/jangle pairs flagged.
\end{verbatim}
}

Within this single, well-constructed instrument, label and content
similarities align; its habitat is cross-instrument screening, where such
fallacies are common \citep{wulff-mata-2025}.
